% LaTeX master Datei(en) zusammengestellt von Falk-Jonatan Strube zur Nutzung an der Hochschule für Technik und Wirtschaft Dresden: https://github.com/genericFJS/htw
\documentclass{scrreprt}
\gdef\pathtomaster{../_LaTeX_master}
%% HTW Dresden Corporate Design
%% Copyright (C) Falk-Jonatan Strube <falkjonatanstrube@gmail.com>, 2017
%%% === Anfang Paket htwscr ===
\makeatletter
\ifcsname KOMAoption\endcsname%
\KOMAoption{footheight}{11mm}	% mehr Platz für Wasserzeichen o.ä. in der Fußzeile
\recalctypearea	% s.o.
\RequirePackage{geometry}
\geometry{top=25mm, bottom=25mm, headsep=10mm, inner=25mm, outer=25mm, footskip=15mm, marginparsep=2mm}
\RequirePackage{helvet}	% Corporate Design schreibt Arial/Helvetica vor
\renewcommand{\familydefault}{\sfdefault}	% s.o.
\RequirePackage{graphicx}	% für Grafiken
\RequirePackage{xcolor}		% für Farben
\definecolor{htworange}{RGB}{249,155,28}	% Corporate Design orange CMYK(0/45/100/0)
\definecolor{htwblue}{RGB}{0,116,188}			% Corporate Design blau CMYK(90/50/0/0)
\definecolor{htwgrey}{RGB}{128,128,128}		% Corporate Design gray 60% Schwarz [Annahme: CMYK(0,0,0,60)]
\RequirePackage{ifthen,xifthen}	% für Bedingungen/Optionen
\RequirePackage{etoolbox}				% s.o.

%%% Format des HTW-Headers (wird hier als Texteelement verwenden) nach dem Style der PowerPoint-Vorlagen. Weiterhin inspiriert durch das Paket tudscr (Corporate Design der TU Dresden)
%%% Dateipfad des Logos abhängig von der Verwendung als Paket/Input
\ifdefined\@isHTWCDpackage%
	\gdef\@pathtoimage{HTW}%
\else%
	\ifdefined\pathtomaster%
		\gdef\@pathtoimage{\pathtomaster /HTW}%
	\else 
		\gdef\@pathtoimage{HTW}%
	\fi%
\fi%
\newcommand{\pathtologo}[1]{\gdef\@pathtoimage{#1}}
%%% Berechnung der Größe des Logos
\newlength{\htwLogoHeight}
\setlength{\htwLogoHeight}{12mm}	% Mindesthöhe des Logos laut Corporate Design: 11mm
\newlength{\@htwLogoWidth}
\AtBeginDocument{
\def\printLogo{\includegraphics[height=\htwLogoHeight]{\@pathtoimage}}
\setlength\@htwLogoWidth{\widthof{\printLogo}}
}
%%% Definition des Header
\newcommand{\@htwheader}{{\noindent%
%%% Header ist abhängig von dem Vorhandensein von department, institute und chair
\newcommand{\hasDepInsCh}{\gdef\@hasDepInsCh{}}%
\ifdefined\@department\hasDepInsCh\fi%
\ifdefined\@institute\hasDepInsCh\fi%
\ifdefined\@chair\hasDepInsCh\fi%
\begin{minipage}[b][\htwLogoHeight]{\dimexpr \textwidth - \@htwLogoWidth - 2mm \relax}% 2mm Abstand zwischen Logo und Linie
{\small{\bfseries%
\ifdefined\@faculty\@faculty{}\fi} %
\strut}%
{\color{htworange}\hrule height 1pt}%
\vspace{.1mm}%
\end{minipage}
\hfill%
\begin{minipage}[b][\htwLogoHeight]{\@htwLogoWidth}
\includegraphics[height=\htwLogoHeight]{\@pathtoimage}
\end{minipage}
\ifdefined\@hasDepInsCh%
\\%
\begin{minipage}{\textwidth}
\small%
\ifdefined\@department\@department{}\fi%
\ifdefined\@institute\ifdefined\@department,~\fi \@institute\fi%
\ifdefined\@chair\ifdefined\@department,~\else\ifdefined\@institute,~\fi\fi\@chair\fi\strut
%{\color{htworange}\hrule height 0.5pt}
\end{minipage}
\fi
}}

%%% Definition der Fußzeilen
\usepackage{scrlayer-scrpage}
\@ifclassloaded{scrreprt}{%
\renewcommand*{\partpagestyle}{scrheadings}
\renewcommand*{\chapterpagestyle}{scrheadings}%
}{}
\pagestyle{scrheadings}
\cfoot{\vfill\pagemark}
\chead{}

%%% Einstellungen
\newcommand{\simpledate}{\gdef\@simpledate{}}		% Datum nicht mit "Eingereicht am:" Vorsatz
\newcommand{\nouppercase}{\gdef\@nouppercase{}}	% Titel/Parts/Chapters usw. nicht in Großbuchstaben

%%% Einiges erst zu Beginn des Dokumentes ausführen, um Einstellungen zu berücksichtigen
\AtBeginDocument{
\renewcommand\SS{SS}	% fix big ß
%%% Sektionen usw. ggf. in Großbuchstaben:
\providecommand{\chapterlinesformat}[3]{}	% compatibility
\providecommand{\chapterlineswithprefixformat}[3]{}	% compatibility
\providecommand{\sectionlinesformat}[4]{}	% compatibility
\ifdefined\@nouppercase\relax\else
\renewcommand\sectionlinesformat[4]{\@hangfrom{\hskip#2 #3}{\MakeUppercase{#4}}}
\@ifclassloaded{scrreprt}{%
\renewcommand\chapterlinesformat[3]{\@hangfrom{#2}{\MakeUppercase{#3}}}
\renewcommand\chapterlineswithprefixformat[3]{\MakeUppercase{#2#3}}%
}{}
\fi
%%% Anpassung der Part-Seite
\renewcommand*{\partheadstartvskip}{\@htwheader\vspace*{30mm}}
\renewcommand*{\raggedpart}{\raggedright}
\renewcommand*{\partformat}{\huge\ifdefined\@nouppercase \partname~\thepart\else\MakeUppercase{\partname~\thepart}\fi.}
\ifdefined\@nouppercase\setkomafont{part}{\Huge}\else\setkomafont{part}{\Huge\MakeUppercase}\fi
}

\newcommand{\bilingual}[3]{\ifdefined #1 \let #1\undefined\fi\newcommand{#1}{\iflanguage{english}{#3}{#2}}}	% Übersetzung von einigen Variablen in mehrere Sprachen

\newcommand{\checkThesis}[2]{\ifthenelse{\equal{#1}{diss}\OR\equal{#1}{doctoral}\OR\equal{#1}{phd}}{
\gdef#2{\dissertationname}}{
\ifthenelse{\equal{#1}{diploma}}{
\gdef#2{\diplomathesisname}}{
\ifthenelse{\equal{#1}{master}}{
\gdef#2{\masterthesisname}}{
\ifthenelse{\equal{#1}{bachelor}}{
\gdef#2{\bachelorthesisname}}{
\ifthenelse{\equal{#1}{student}}{
\gdef#2{\studentthesisname}}{
\ifthenelse{\equal{#1}{evidence}}{
\gdef#2{\studentresearchname}}{
\ifthenelse{\equal{#1}{project}}{
\gdef#2{\projectpapername}}{
\ifthenelse{\equal{#1}{seminar}}{
\gdef#2{\seminarpapername}}{
\ifthenelse{\equal{#1}{term}}{
\gdef#2{\termpapername}}{
\ifthenelse{\equal{#1}{research}}{
\gdef#2{\researchname}}{
\ifthenelse{\equal{#1}{log}}{
\gdef#2{\logname}}{
\ifthenelse{\equal{#1}{report}}{
\gdef#2{\reportname}}{
\ifthenelse{\equal{#1}{internship}}{
\gdef#2{\internshipname}}{
\ifthenelse{\equal{#1}{lecture}\OR\equal{#1}{lesson}\OR\equal{#1}{pract}}{
\bilingual{\professorname}{Vorlesung von}{Lecture by}
\simpledate
\gdef\@islecture{\relax}
\ifthenelse{\equal{#1}{lecture}}{\gdef#2{\lecturename}}{
\ifthenelse{\equal{#1}{lesson}}{\gdef#2{\lessonname}}{
\gdef#2{\practname}
}}}{
\gdef#2{#1}}}}}}}}}}}}}}}}

%%% Textfragmente der Titelseite in deutsch und englisch
\bilingual{\dateofbirthname}{Geboren am:}{Born on:}
\bilingual{\placeofbirthname}{in}{in}
\bilingual{\coursename}{Studiengang:}{Course:}
\bilingual{\disciplinename}{Studienrichtung:}{Discipline:}
\bilingual{\matriculationnumbername}{Matrikelnummer:}{Matriculation number:}
\bilingual{\matriculationyearname}{Immatrikulationsjahr}{Matriculation year:}
\bilingual{\graduationname}{zur Erlangung des akademischen Grades}{to achieve the academic degree}
\bilingual{\refereename}{Gutachter}{Referee}
\bilingual{\advisorname}{Fachreferent}{Advisor}
\bilingual{\supervisorname}{Betreuer}{Supervisor}
\bilingual{\professorname}{Betreuender Hochschullehrer}{Supervising professor}
\bilingual{\datename}{Eingereicht am:}{Submitted on:}
%% Thesis Typen (inspiriert von dem Paket tudscr):
\bilingual{\dissertationname}{Dissertation}{Dissertation} % diss, doctoral, phd 
\bilingual{\diplomathesisname}{Diplomarbeit}{Diploma Thesis} % diploma 
\bilingual{\masterthesisname}{Master-Arbeit}{Master Thesis} % master 
\bilingual{\bachelorthesisname}{Bachelor-Arbeit}{Bachelor Thesis} % bachelor 
\bilingual{\studentthesisname}{Studienarbeit}{Student Thesis} % student 
\bilingual{\studentresearchname}{Großer Beleg}{Student Research Project} % evidence 
\bilingual{\projectpapername}{Projektarbeit}{Project Paper} % project
\bilingual{\seminarpapername}{Seminararbeit}{Seminar Paper} % seminar
\bilingual{\termpapername}{Hausarbeit}{Term Paper} % term
\bilingual{\researchname}{Forschungsbericht}{Research Report} % research 
\bilingual{\logname}{Protokoll}{Log} % log 
\bilingual{\reportname}{Bericht}{Report} % report
\bilingual{\internshipname}{Praktikumsbericht}{Internship Report} % internship 
\bilingual{\lecturename}{Vorlesungsmitschrift}{Lecture Notes} % lecture 
\bilingual{\lectureauthorname}{Mitschrift von}{Notes by} % lecture  
\bilingual{\lessonname}{\"Ubungsmitschrift}{Lesson Notes} % lecture  
\bilingual{\practname}{Praktikumsmitschrift}{Pracital Course Notes} % lecture 
%%% Variablen für die Titelseite (inspiriert von dem Paket tudscr):
\renewcommand{\subject}[1]{\checkThesis{#1}{\@subject}}
\newcommand{\faculty}[1]{\gdef\@faculty{#1}}
\newcommand{\department}[1]{\gdef\@department{#1}}
\newcommand{\institute}[1]{\gdef\@institute{#1}}
\newcommand{\chair}[1]{\gdef\@chair{#1}}
\newcommand{\authormore}[1]{\gdef\@authormore{#1}}
\newcommand{\dateofbirth}[1]{\gdef\@dateofbirth{#1}}
\newcommand{\placeofbirth}[1]{\gdef\@placeofbirth{#1}}
\newcommand{\course}[1]{\gdef\@course{#1}}
\newcommand{\discipline}[1]{\gdef\@discipline{#1}}
\newcommand{\matriculationnumber}[1]{\gdef\@matriculationnumber{#1}}
\newcommand{\matriculationyear}[1]{\gdef\@matriculationyear{#1}}
\newcommand{\thesis}[1]{\checkThesis{#1}{\@thesis}}
\newcommand{\graduation}[2][]{\gdef\@graduation{#2}\gdef\@graduationshort{#1}}
\newcommand{\supervisor}[1]{\gdef\@supervisor{#1}}
\newcommand{\referee}[1]{\gdef\@referee{#1}}
\newcommand{\advisor}[1]{\gdef\@advisor{#1}}
\newcommand{\professor}[1]{\gdef\@professor{#1}}
%%% Titelseite (inspiriert von dem Paket tudscr):
\renewcommand*{\maketitle}{%
\begin{titlepage}%
\begin{flushleft}
\renewcommand{\thefootnote}{$\star$}%
%%% Kopfzeile HTW-Style:
\@htwheader\\[13mm]
%%% Die Zeilen für titlehead, subject, title, subtitle und author sind immer vorhanden (spacing)
{\noindent\@titlehead{}\strut}\\[3mm]%
{\Large\@subject{}\strut}\\[8mm]%
{\Huge\bfseries%
\ifdefined\@nouppercase\@title{}\else\MakeUppercase{\@title{}}\fi%
\strut}\\
{\Large\bfseries%
\ifdefined\@nouppercase\@subtitle{}\else\MakeUppercase{\@subtitle{}}\fi%
\strut}\\[11mm]
\ifdefined\@islecture\strut\else{\Large\@author{}\strut}\fi % Autor wird bei Vorlesungsmitschrift unten angezeigt
%%% Weitere Felder werden nur angezeigt, wenn sie definiert wurden
\ifdefined\@authormore\ifdefined\@islecture\else\\\fi{\@authormore{}\strut}\fi
\ifdefined\@dateofbirth\\\dateofbirthname{} \@dateofbirth{} \ifdefined\@placeofbirth\placeofbirthname{} \@placeofbirth{}\strut\fi\fi
\ifdefined\@course\\\coursename{} \@course{}\strut\fi
\ifdefined\@discipline\\\disciplinename{} \@discipline{}\strut\fi
\ifdefined\@matriculationnumber\\\matriculationnumbername{} \@matriculationnumber{}\strut\fi
\ifdefined\@matriculationyear\\\matriculationyearname{} \@matriculationyear{}\strut\fi
\ifdefined\@thesis\\[10mm]{\LARGE\bfseries
\ifdefined\@nouppercase\@thesis{}\else\MakeUppercase{\@thesis{}}\fi
\strut}
\ifdefined\@graduation\\[2mm]{\graduationname{}\strut}
\\[5mm]{\LARGE\bfseries
\ifdefined\@nouppercase\relax\else\MakeUppercase{\fi%
\@graduation{} \ifdefined\@graduationshort(\@graduationshort{})\fi
\ifdefined\@nouppercase\relax\else}\fi\strut}\fi\fi
\vfill\noindent
\ifdefined\@islecture\\[1mm]{\lectureauthorname{}\strut}\\[.7mm]{\Large\@author{}\strut}\fi
\ifdefined\@referee{\refereename{}\strut}\\[.7mm]{\Large\@referee{}\strut}\fi
\ifdefined\@advisor\\[1mm]{\advisorname{}\strut}\\[.7mm]{\Large\@advisor{}\strut}\fi
\ifdefined\@supervisor\\[1mm]{\supervisorname{}\strut}\\[.7mm]{\Large\@supervisor{}\strut}\fi
\ifdefined\@professor\\[1mm]{\professorname{}\strut}\\[.7mm]{\Large\@professor{}\strut}\fi
\ifdefempty{\@date}{}{\ifdefined\@simpledate\\[6mm]\relax\else\\[3mm]\datename{}~\fi\@date{}\strut}
\ifdefempty{\@publishers}{}{\\[4.5mm]{\Large\@publishers{}\strut}}%
\renewcommand\footnoterule{}
\@thanks\global\let\@thanks\@empty
\end{flushleft}
\end{titlepage}
\setcounter{footnote}{0}
}
\else%
\PackageError{htwcd}{htwcd needs a KOMA class!}%
\fi%
\makeatother
%%% === Ende Paket htwscr ===
%-- Nummerierung bis Subsubsection für Report
\setcounter{secnumdepth}{3}	% zählt auch subsubsection
\setcounter{tocdepth}{3}	% Inhaltsverzeichnis bis in subsubsection

% LuaLaTeX Unterscheidung
\usepackage{ifluatex}

% Input inkl. Umlaute, Silbentrennung
\ifluatex
	\usepackage{fontspec}
\else
	\usepackage[utf8]{inputenc}
	\usepackage[T1]{fontenc}
\fi
\usepackage[english,ngerman]{babel}
\usepackage{csquotes}	% Anführungszeichen
\usepackage{marvosym}
\usepackage{eurosym}

% Style-Aufhübschung
\usepackage{soul, color}	% Kapitälchen, Unterstrichen, Durchgestrichen usw. im Text
%\usepackage{titleref}
\usepackage[breakwords, fit]{truncate}	% Abschneiden von Sätzen
\renewcommand{\TruncateMarker}{\,…}

% Mathe usw.
\usepackage{amssymb}
\usepackage{amsthm}
\usepackage[fleqn,intlimits]{amsmath}	% fleqn: align-Umgebung rechtsbündig; intlimits: Integralgrenzen immer ober-/unterhalb
%\usepackage{mathtools} % u.a. schönere underbraces
\usepackage{xcolor}
\usepackage{esint}	% Schönere Integrale, \oiint vorhanden
\everymath=\expandafter{\the\everymath\displaystyle}	% Mathe Inhalte werden weniger verkleinert
\usepackage{wasysym}	% mehr Symbole, bspw \lightning
%\renewcommand{\int}{\int\limits}
%\usepackage{xfrac}	% mehr fracs: sfrac{}{}
\let\oldemptyset\emptyset	% schöneres emptyset
\let\emptyset\varnothing
%\RequirePackage{mathabx}	% mehr Symbole
\mathchardef\mhyphen="2D	% Hyphen in Math
\usepackage{mathtools}		% zum Beispiel {pmatrix}[r] für Formatierung der Matrix

% tikz usw.
\usepackage{tikz}
\usepackage{pgfplots}
\pgfplotsset{compat=1.11}	% Umgeht Fehlermeldung
\usetikzlibrary{graphs}
%\usetikzlibrary{through}	% ???
\usetikzlibrary{arrows}
\usetikzlibrary{arrows.meta}	% Pfeile verändern / vergrößern: \draw[-{>[scale=1.5]}] (-3,5) -> (-3,3);
\usetikzlibrary{automata,positioning} % Zeilenumbruch im Node node[align=center] {Text\\nächste Zeile} automata für Graphen
\usetikzlibrary{matrix}
\usetikzlibrary{patterns}	% Schraffierte Füllung
\usetikzlibrary{shapes.geometric}	% Polygon usw.
\tikzstyle{reverseclip}=[insert path={	% Inverser Clip \clip
	(current page.north east) --
	(current page.south east) --
	(current page.south west) --
	(current page.north west) --
	(current page.north east)}
% Nutzen: 
%\begin{tikzpicture}[remember picture]
%\begin{scope}
%\begin{pgfinterruptboundingbox}
%\draw [clip] DIE FLÄCHE, IN DER OBJEKT NICHT ERSCHEINEN SOLL [reverseclip];
%\end{pgfinterruptboundingbox}
%\draw DAS OBJEKT;
%\end{scope}
%\end{tikzpicture}
]	% Achtung: dafür muss doppelt kompliert werden!
\usepackage{graphpap}	% Grid für Graphen
\tikzset{every state/.style={inner sep=2pt, minimum size=2em}}
\usetikzlibrary{mindmap, backgrounds}
%\usepackage{tikz-uml}	% braucht Dateien: http://perso.ensta-paristech.fr/~kielbasi/tikzuml/

% Tabular
\usepackage{longtable}	% Große Tabellen über mehrere Seiten
\usepackage{multirow}	% Multirow/-column: \multirow{2[Anzahl der Zeilen]}{*[Format]}{Test[Inhalt]} oder \multicolumn{7[Anzahl der Reihen]}{|c|[Format]}{Test2[Inhalt]}
\renewcommand{\arraystretch}{1.3} % Tabellenlinien nicht zu dicht
\usepackage{colortbl}
\arrayrulecolor{gray}	% heller Tabellenlinien
\usepackage{array}	% für folgende 3 Zeilen (für Spalten fester breite mit entsprechender Ausrichtung):
\newcolumntype{L}[1]{>{\raggedright\let\newline\\\arraybackslash\hspace{0pt}}m{\dimexpr#1\columnwidth-2\tabcolsep-1.5\arrayrulewidth}}
\newcolumntype{C}[1]{>{\centering\let\newline\\\arraybackslash\hspace{0pt}}m{\dimexpr#1\columnwidth-2\tabcolsep-1.5\arrayrulewidth}}
\newcolumntype{R}[1]{>{\raggedleft\let\newline\\\arraybackslash\hspace{0pt}}m{\dimexpr#1\columnwidth-2\tabcolsep-1.5\arrayrulewidth}}
\usepackage{booktabs}	% Kompatibilität für array und \hrule
\usepackage{caption}	% Um auch unbeschriftete Captions mit \caption* zu machen

% Nützliches
\usepackage{verbatim}	% u.a. zum auskommentieren via \begin{comment} \end{comment}
\usepackage{tabto}	% Tabs: /tab zum nächsten Tab oder /tabto{.5 \CurrentLineWidth} zur Stelle in der Linie
\NumTabs{6}	% Anzahl von Tabs pro Zeile zum springen
\usepackage{listings} % Source-Code mit Tabs
\usepackage{lstautogobble} 
\usepackage{enumitem}	% Anpassung der enumerates
%\setlist[enumerate,1]{label=(\arabic*)}	% global andere Enum-Items
\renewcommand{\labelitemiii}{$\scriptscriptstyle ^\blacklozenge$} % global andere 3. Item-Aufzählungszeichen
\usepackage{letltxmacro} % neue Definiton von Grundbefehlen
% Nutzen:
%\LetLtxMacro{\oldemph}{\emph}
%\renewcommand{\emph}[1]{\oldemph{#1}}
\RequirePackage{xpatch}	% ua. Konkatenieren von Strings/Variablen (etoolbox)
\usepackage{xstring}	% String Operationen
\usepackage{minibox}	% Minibox anstatt \fbox{} für Boxen mit (manuellen) Zeilenumbruch
\usepackage{framed}		% \begin{framed} statt minibox/fbox mit automatischem Zeilenumbruch
\usepackage{microtype}	% Verfeinerte Hyphenation
\usepackage{import}		
\usepackage{standalone}	% Zusammen mit import: Einfügen von separaten Latex-Dateien (die schon header haben)

% Einrichtung von lst
%\renewcommand*{\ttdefault}{qcr}
\lstset{
basicstyle=\ttfamily, 
%basicstyle=\footnotesize\sffamily\color{black},
%mathescape=true, 
%escapeinside=^^, 
autogobble, 
tabsize=2,
frame=single,
rulecolor=\color{lightgray},
numbers=left,
numbersep=5pt,
numberstyle=\tiny\color{gray},
commentstyle=\color{gray},
keywordstyle=\color{htwblue},
stringstyle=\color{htworange},
morecomment=[l][\color{magenta}]{\#}
showspaces=false,
showstringspaces=false,
breaklines=true,
literate=%
    {Ö}{{\"O}}1
    {Ä}{{\"A}}1
    {Ü}{{\"U}}1
    {ß}{{\ss}}1
    {ü}{{\"u}}1
    {ä}{{\"a}}1
    {ö}{{\"o}}1
    {„}{{\glqq}}1
    {“}{{\grqq}}1
    {~}{{\textasciitilde}}1
}
\usepackage{scrhack} % Fehler umgehen
\def\ContinueLineNumber{\lstset{firstnumber=last}} % vor lstlisting. Zum wechsel zum nicht-kontinuierlichen muss wieder \StartLineAt1 eingegeben werden
\def\StartLineAt#1{\lstset{firstnumber=#1}} % vor lstlisting \StartLineAt30 eingeben, um bei Zeile 30 zu starten
\let\numberLineAt\StartLineAt

% BibTeX
\usepackage[
backend=bibtex,
bibencoding=ascii,
style=authoryear, citestyle=authoryear-ibid,
%doi=false,
%isbn=false,
%url=false
]{biblatex}	% BibTeX
\usepackage{makeidx}
%\makeglossary
%\makeindex
\usepackage{url}	% Darstellung von URLs

% Grafiken
\usepackage{graphicx}
\usepackage{epstopdf}	% eps-Vektorgrafiken einfügen
\usepackage{transparent}	% transparent nutzen: {\transparent{0.4} ...}
%\epstopdfsetup{outdir=\customDir}
% Prüft, ob Grafik existiert (mit \ifvalidimage{}{}) [Quelle: https://tex.stackexchange.com/a/99176]:
\makeatletter
\newif\ifgraphicexist
\catcode`\*=11
\newcommand\ifvalidimage[1]{%
    \begingroup
    \global\graphicexisttrue
    \let\input@path\Ginput@path
    \filename@parse{#1}%
    \ifx\filename@ext\relax
    \@for\Gin@temp:=\Gin@extensions\do{%
        \ifx\Gin@ext\relax
        \Gin@getbase\Gin@temp
        \fi}%
    \else
    \Gin@getbase{\Gin@sepdefault\filename@ext}%
    \ifx\Gin@ext\relax
    \global\graphicexistfalse
    \def\Gin@base{\filename@area\filename@base}%
    \edef\Gin@ext{\Gin@sepdefault\filename@ext}%
    \fi
    \fi
    \ifx\Gin@ext\relax
    \global\graphicexistfalse
    \else 
    \@ifundefined{Gin@rule@\Gin@ext}%
    {\global\graphicexistfalse}%
    {}%
    \fi  
    \ifx\Gin@ext\relax 
    \gdef\imageextension{unknown}%
    \else
    \xdef\imageextension{\Gin@ext}%
    \fi 
    \endgroup 
    \ifgraphicexist
    \expandafter \@firstoftwo
    \else
    \expandafter \@secondoftwo
    \fi 
} 
\catcode`\*=12
\makeatother
\LetLtxMacro{\forceincludegraphics}{\includegraphics}	% neuer Befehl für includegraphics
\renewcommand{\includegraphics}[2][]{	% altes includegraphics neu definieren, damit es auch nicht vorhandene einfügt
\ifvalidimage{#2}{
\forceincludegraphics[#1]{#2}
}{
\message{Achtung: Grafik wurde nicht gefunden: '#2'}
\minibox[frame]{
\textbf{\StrSubstitute{#2}{_}{\_}}%  \ifthenelse{\isempty{#1}}{}{\\\textit{#1}}
}}}

% pdf-Setup
\usepackage{pdfpages}
\usepackage[bookmarks,%
bookmarksopen=false,% Klappt die Bookmarks in Acrobat aus
colorlinks=true,%
linkcolor=black,%
citecolor=black,%
urlcolor=htworange,%
]{hyperref}

% Nach Dokumentenbeginn ausführen:
\makeatletter
\AtBeginDocument{	
	% pdf Informationen
	\hypersetup{
		pdfauthor={\@author},
		pdftitle={\@title},
		pdfproducer={LaTeX Layout erstellt von Falk-Jonatan Strube}
	}
	% signierte Fußzeile
	%\providecommand{\ofootname}{Mitschrift von}
	\providecommand{\pathtomaster}{../}
	\ofoot{%
	%\ifdefined\ofoottext
	%\else
	%\begin{minipage}{40mm}
	%\begin{flushright}
	%\textcolor{gray}{
	%\ofootname\\
	%Falk-Jonatan Strube\vspace*{-1mm}}
	%\end{flushright}
	%\end{minipage}
	%\fi%
	%\hspace*{1mm}%
	\begin{minipage}{10 mm}
	\includegraphics[height=9mm]{\pathtomaster/FJS-logo_gray}
	\end{minipage}%
	}
}
\makeatother



%%% ==================================================
%%% ===================== MACROS =====================
%%% ==================================================
% Horizontale Linie:
\newcommand{\HRule}[1][\medskipamount]{\par
  \vspace*{\dimexpr-\parskip-\baselineskip+#1}
  \noindent\rule[0.2ex]{\linewidth}{0.2mm}\par
  \vspace*{\dimexpr-\parskip-.5\baselineskip+#1}}
% Gestrichelte horizontale Linie:
\RequirePackage{dashrule}
\newcommand{\HDRule}[1][\medskipamount]{\par
  \vspace*{\dimexpr-\parskip-\baselineskip+#1}
  \noindent\hdashrule[0.2ex]{\linewidth}{0.2mm}{1mm} \par
  \vspace*{\dimexpr-\parskip-.5\baselineskip+#1}}
% Mathe in Anführungszeichen:
\newsavebox{\mathbox}\newsavebox{\mathquote}
\makeatletter
\newcommand{\mq}[1]{% \mathquotes{<stuff>}
  \savebox{\mathquote}{\text{"}}% Save quotes
  \savebox{\mathbox}{$\displaystyle #1$}% Save <stuff>
  \raisebox{\dimexpr\ht\mathbox-\ht\mathquote\relax}{"}#1\raisebox{\dimexpr\ht\mathbox-\ht\mathquote\relax}{''}
}
\makeatother

% Paragraph mit Zähler (Section-Weise)
\newcounter{cparagraphC}
\newcommand{\cparagraph}[1]{
\stepcounter{cparagraphC}
\paragraph{\thesection{}-\thecparagraphC{} #1}
%\addcontentsline{toc}{subsubsection}{\thesection{}-\thecparagraphC{} #1}
\label{\thesection-\thecparagraphC}
}
\makeatletter
\@addtoreset{cparagraphC}{section}
\makeatother


% (Vorlesungs-)Folien einbinden:
\makeatletter
% Folien-Pfad (inkl. Dateiname ohne Endung und ggf. ohne Nummerierung)
\newcommand{\@slidePath}{}
\newcommand{\slidePath}[1]{\renewcommand{\@slidePath}{#1}}
% Folien Eigenschaften
\newcommand{\@slideScale}{0.5}
\newcommand{\slideScale}[1]{\renewcommand{\@slideScale}{#1}}
\AtBeginDocument{
% Folien von einer Datei skaliert
\newcommand{\slide}[2][\@slideScale]{\slides[#1]{}{#2}}
\newcommand{\slideTrim}[6][\@slideScale]{\slides[#1 , clip,  trim = #5cm #4cm #6cm #3cm]{}{#2}}
% Folien von mehreren nummerierten Dateien skaliert
\newcommand{\slides}[3][\@slideScale]{\begin{center}
\includegraphics[page=#3, scale=#1]{\@slidePath #2.pdf}
\end{center}}
}
\makeatother

% \emph{} anders definieren
\makeatletter
\DeclareRobustCommand{\em}{%
  \@nomath\em \if b\expandafter\@car\f@series\@nil
  \normalfont \else \scshape \fi}
\makeatother

% unwichtiges
\newcommand{\unimptnt}[1]{{\transparent{0.5}#1}}

% alph. enumerate
\newenvironment{anumerate}{\begin{enumerate}[label=(\alph*)]}{\end{enumerate}} % Alphabetische Aufzählung

% Hanging parameters
\newcommand{\hangpara}[1]{\par\noindent\hangindent+2em\hangafter=1 #1\par\noindent}

%% EINFACHE BEFEHLE

% Abkürzungen Mathe
\newcommand{\EE}{\mathbb{E}}
\newcommand{\QQ}{\mathbb{Q}}
\newcommand{\RR}{\mathbb{R}}
\newcommand{\CC}{\mathbb{C}}
\newcommand{\NN}{\mathbb{N}}
\newcommand{\ZZ}{\mathbb{Z}}
\newcommand{\PP}{\mathbb{P}}
\renewcommand{\SS}{\mathbb{S}}
\newcommand{\cA}{\mathcal{A}}
\newcommand{\cB}{\mathcal{B}}
\newcommand{\cC}{\mathcal{C}}
\newcommand{\cD}{\mathcal{D}}
\newcommand{\cE}{\mathcal{E}}
\newcommand{\cF}{\mathcal{F}}
\newcommand{\cG}{\mathcal{G}}
\newcommand{\cH}{\mathcal{H}}
\newcommand{\cI}{\mathcal{I}}
\newcommand{\cJ}{\mathcal{J}}
\newcommand{\cM}{\mathcal{M}}
\newcommand{\cN}{\mathcal{N}}
\newcommand{\cP}{\mathcal{P}}
\newcommand{\cR}{\mathcal{R}}
\newcommand{\cS}{\mathcal{S}}
\newcommand{\cZ}{\mathcal{Z}}
\newcommand{\cL}{\mathcal{L}}
\newcommand{\cT}{\mathcal{T}}
\newcommand{\cU}{\mathcal{U}}
\newcommand{\cX}{\mathcal{X}}
\newcommand{\cV}{\mathcal{V}}
\renewcommand{\phi}{\varphi}
\renewcommand{\epsilon}{\varepsilon}
\renewcommand{\theta}{\vartheta}

% Verschiedene als Mathe-Operatoren
\DeclareMathOperator{\arccot}{arccot}
\DeclareMathOperator{\arccosh}{arccosh}
\DeclareMathOperator{\arcsinh}{arcsinh}
\DeclareMathOperator{\arctanh}{arctanh}
\DeclareMathOperator{\arccoth}{arccoth} 
\DeclareMathOperator{\var}{Var} % Varianz 
\DeclareMathOperator{\cov}{Cov} % Co-Varianz 

% Farbdefinitionen
\definecolor{red}{RGB}{180,0,0}
\definecolor{green}{RGB}{75,160,0}
\definecolor{blue}{RGB}{0,75,200}
\definecolor{orange}{RGB}{255,128,0}
\definecolor{yellow}{RGB}{255,245,0}
\definecolor{purple}{RGB}{75,0,160}
\definecolor{cyan}{RGB}{0,160,160}
\definecolor{brown}{RGB}{120,60,10}

\definecolor{itteny}{RGB}{244,229,0}
\definecolor{ittenyo}{RGB}{253,198,11}
\definecolor{itteno}{RGB}{241,142,28}
\definecolor{ittenor}{RGB}{234,98,31}
\definecolor{ittenr}{RGB}{227,35,34}
\definecolor{ittenrp}{RGB}{196,3,125}
\definecolor{ittenp}{RGB}{109,57,139}
\definecolor{ittenpb}{RGB}{68,78,153}
\definecolor{ittenb}{RGB}{42,113,176}
\definecolor{ittenbg}{RGB}{6,150,187}
\definecolor{itteng}{RGB}{0,142,91}
\definecolor{ittengy}{RGB}{140,187,38}

% Textfarbe ändern
\newcommand{\tred}[1]{\textcolor{red}{#1}}
\newcommand{\tgreen}[1]{\textcolor{green}{#1}}
\newcommand{\tblue}[1]{\textcolor{blue}{#1}}
\newcommand{\torange}[1]{\textcolor{orange}{#1}}
\newcommand{\tyellow}[1]{\textcolor{yellow}{#1}}
\newcommand{\tpurple}[1]{\textcolor{purple}{#1}}
\newcommand{\tcyan}[1]{\textcolor{cyan}{#1}}
\newcommand{\tbrown}[1]{\textcolor{brown}{#1}}

% Umstellen der Tabellen Definition
\newcommand{\mpb}[1][.3]{\begin{minipage}{#1\textwidth}\vspace*{3pt}}
\newcommand{\mpe}{\vspace*{3pt}\end{minipage}}

\newcommand{\resultul}[1]{\underline{\underline{#1}}}
\newcommand{\parskp}{$ $\\}	% new line after paragraph
\newcommand{\parnl}{\unskip\parfillskip 0pt \par} % justified new line (brace line!)

% mathematische Abkürzungen
\newcommand{\corr}{\;\widehat{=}\;}
\newcommand{\mdeg}{^{\circ}}
\newcommand{\nok}[2]{\binom{#1}{#2}}	% n über k BESSER: \binom{n}{k}
\newcommand{\mtr}[1]{\begin{pmatrix}#1\end{pmatrix}}	% Matrix
\newcommand{\dtr}[1]{\begin{vmatrix}#1\end{vmatrix}}	% Determinante (Betragsmatrix)
\LetLtxMacro{\originalVec}{\vec}
\renewcommand{\vec}[1]{\underline{#1}}	% Vektorschreibweise
\newcommand{\imptnt}[1]{\colorbox{red!30}{#1}}	% Wichtiges
\newcommand{\intd}[1]{\,\mathrm{d}#1}
\newcommand{\diffd}[1]{\mathrm{d}#1}
% für Module-Rechnung: \pmod{}
\newcommand{\unit}[1]{\,\mathrm{#1}}
%\LetLtxMacro{\ntilde}{\tilde}
%\renewcommand{\tilde}{\widetilde}

% weiter Abkürzungen
\newcommand{\gdw}{genau dann wenn}
\newcommand{\lecdate}[1]{\marginpar{\small\textcolor{gray}{Vorlesung\newline #1}}}

\faculty{Fakultät Informatik/Mathematik}
%\department{Einrichtung}
%\institute{Institut}
\chair{Lehrstuhl für Künstliche Intelligenz / Theoretische Informatik}
\subject{lecture}
\title{Künstliche Intelligenz}
%\subtitle{Untertitel}
\author{Falk-Jonatan Strube}
%\thesis{Typisierung Thesis}
\professor{Prof. Dr. Boris Hollas}
%\date{Datum}

\begin{document}

\maketitle
\tableofcontents

\chapter*{Vorbemerkungen}
\section*{Inhalt}
\begin{enumerate}
\item Prädikatenlogik
\item Prolog-Programmmierung
\item Logisches Schließen
\item Sprachverstehen
\item Problemlösen durch Suche
\item Bayes'sche Netze
\end{enumerate}

\section*{Prüfungsvorleistung}
Alle 1-2 Wochen werden Aufgaben gestellt, die innerhalb von 2 Wochen gelöst und vorgeführt werden müssen. Mindestens 4 dieser Aufgaben müssen erfolgreich gelöst werden.

\section*{Prüfung}
\begin{itemize}
\item selbstbeschriebenes Blatt
\item Taschenrechner
\end{itemize}

\section*{Literatur}
\begin{itemize}
\item Ertl: Grundkurs Künstliche Intelligenz
\item Russell/Norvig: Künstliche Intelligenz
\end{itemize}

\chapter{Prädikatenlogik}
\lecdate{20.03.2017}
Problem: Aussagenlogik ist wenig mächtig.\\
Aussagen, die sich nicht formulieren lassen:
\begin{itemize}
\item „Alle Vögel können fliegen.“\\
(Nur formulierbar, wenn alle Vögel einzeln als Atomformeln beschrieben werden. Diese Formel der Aussagenlogik wäre dann allerdings sehr lang.)
\item „Wenn X eine Katze ist, dann ist X ein Haustier.“\\
(In der Aussagenlogik sind keine Variablen möglich.)
\item „Für jedes Land gibt es eine Hauptstadt.“
\end{itemize}

\section{Syntax der Prädikatenlogik}
\cparagraph{Definition}
Sei $V$ eine Menge von Variablen, $K$ eine Menge von Konstante und $F$ eine Menge von Funktionssymbolen.
\begin{itemize}
\item Dann sind alle Variablen $V$ und Konstanten in $K$ Terme.
\item Wenn $t_1, \dots, t_n$ Terme sind und $f$ ein $n$-stelliges Funktionssymbol ist, dann ist auch $f(t_1, \dots, t_n)$ ein Term. 
\end{itemize}

\cparagraph{Beispiel}
$V=\{x,y,z\},\; K=\{a,b,c\},\; F=\{+,*\}$. Terme sind $x,a,x+a,x*(x+1)$.

\cparagraph{Definition}
Sei eine Menge von Prädikatensymbolen gegeben. Die Formeln der Prädikatenlogik sind induktiv definiert:
\begin{itemize}
\item Wenn $t_1, \dots, t_n$ Term und $P$ ein Prädikatensymbol der Stelligkeit $n$ ist,\\ dann ist $P(t_1, \dots, t_n)$ eine Formel.
\item Sind $F,G$ Formeln, dann auch $F\wedge G, \; F\vee G,\; \neg F,\; (F)$.
\item Wenn $x$ eine Variable und $F$ eine Formel ist, dann sind auch $\forall x \,F, \; \exists x\,F$ Formeln.
\end{itemize}

\cparagraph{Bemerkung} Wie in der Aussagenlogik definiert sich $\to, \; \leftrightarrow$.\\
Vorrang der Operatoren: $\neg;\; \forall, \; \exists;\; \wedge, \; \vee;\; \to,\; \leftrightarrow$

\cparagraph{Beispiel}
\begin{itemize}
\item $\forall x\, (vogel(x) \to fliegt(x))$\\
\begin{tikzpicture}
\node (v1) at (-0.5,2) {$\forall x$};
\node (v2) at (-0.5,1) {$\to$};
\node (v3) at (-1.5,-0.5) {$vogel(x)$};
\node (v4) at (0.5,-0.5) {$fliegt(x)$};
\draw (v1) -- (v2);
\draw (v2) -- (v3);
\draw (v2) -- (v4);
\end{tikzpicture}
\item $\forall x\, (katze(x) \to haustier(x))$
\item $\forall x\, (land(x) \to \exists y\, hauptstadt (x,y))$\\
\begin{tikzpicture}
\node (v1) at (-0.5,2) {$\forall x$};
\node (v2) at (-0.5,1) {$\to$};
\node (v3) at (-1.5,-0.5) {$land(x)$};
\node (v4) at (0.5,-0.5) {$\exists y$};
\draw (v1) -- (v2);
\draw (v2) -- (v3);
\draw (v2) -- (v4);
\node (v5) at (0.5,-2) {$hauptstadt(x,y)$};
\draw (v4) -- (v5);
\end{tikzpicture}
\item $katz(feli)$ mit $feli \in K$
\end{itemize}

\section{Semantik der Prädikatenlogik (informal)}
Wir betrachten zunächst eine (unwichtige) Menge $U$ (Universum), die alle zu betrachtenden Objekte (Vögel, Katzen, Länder usw.) enthält. Davon betrachten wir Teilmengen, z.B. die Menge aller Vögel, Länder usw. sowie mehrstellige Relationen auf $U$ (z.B. Hauptstadt, verheiratet).

Beispiel: $Hauptstadt=\{(Berlin,\;Deutschland), (Paris, \; Frankreich)\} \subseteq \underbrace{U \times U}_{oder:\; Stadt\times Land}$\\
Den Prädikatensymbolen müssen Prädikate (Relationen) zugeordnet werden.

Beispiel: Ordnen dem Prädikatensymbol $vogel$ das Prädikat „Menge aller Vögel“ zu.\\
Dadurch ergibt sich der Wahrheitswert einer Formel: 
\begin{itemize}
\item Die Formel $P(t_1, \dots, t_n)$ ist wahr \gdw{} $(t_1, \dots, t_n) \in P$.
\item Die Wahrheit von $F\wedge G, \; F\vee G,\; \neg F$ ist wie in der Aussagenlogik definiert.
\item $\exists x \, F$ ist wahr \gdw{} es ein $x$ aus dem Universum gibt, mit dem $F$ wahr ist.
\item $\forall x \, F$ ist wahr \gdw{} $F$ für alle $x$ aus dem Universum wahr ist.
\end{itemize}
\cparagraph{Beispiel}
Sei $vogel=\{Amsel, Drossel, Fink, Star\}$\\
und $fliegt=vogel \cup \{Maikaefer, A380\}$
\begin{itemize}
\item $vogel(Amsel)$ ist wahr.
\item $vogel(Maikaefer)$ ist falsch.
\item $\exists x \, vogel(x)$ ist wahr.
\item $\forall x \, vogel(x)$ ist falsch.
\item $\forall x \,(vogel(x) \to fliegt(x))$ ist wahr.
\end{itemize}

\cparagraph{Beispiel} Sei $land=\{Deutschland, Frankreich\}$ und \\
$hauptstadt = \{ (Berlin, Deutschland), (Paris, Frankreich) \}$.\\
Für jedes Land gibt es eine Hauptstadt:
$$\forall x \,(land(x)\to \exists y \, hauptstadt(y,x))$$
\cparagraph{Beispiel} Seien die Relationen $Katze$ (Menge aller Hauskatzen) sowie $Haustier$ (Menge aller Haustiere) gegeben.\\
Jede Katze ist ein Haustier:
$$\forall x \,(Katze (x) \to Haustier (x))$$
\cparagraph{Beispiel}
Seien $Mensch, \; Mann, \; Frau$ die Menge aller Menschen, Männer und Frauen.
\begin{itemize}
\item Jeder Mensch ist Mann oder Frau:
$$\forall x \, (Mensch(x) \to Mann(x) \vee Frau (x))$$
Anmerkung: wichtige Unterscheidung:\\
$Mensch(x) \to Mann(x) \vee Frau (x)$ liefert falsch zurück, wenn $x$ bspw. ein Stuhl -- das ist logisch richtig und \\
$Mensch(x) \wedge( Mann(x) \vee Frau (x))$ sagt, dass jedes $x$ entweder Mann oder Frau sein muss… ein Stuhl bspw. auch -- das ist falsch!
\item Jeder Mann und jede Frau ist ein Mensch:
$$\forall x \, (Mann(x) \vee Frau(x) \to Mensch(x))$$
Anmerkung: Obwohl im Satz „und“ steht, steht in der Prädikatenlogik das „oder“. Das „und“ ist in diesem Sinne gemeint:\\
$\forall x \, (Mann(x)\to Mensch(x)) \wedge \forall x\, (Frau(x) \to Mensch(x))$
\item Kein Mensch ist gleichzeitig Mann und Frau:
$$\forall x \, (Mensch(x) \to  \neg (Mann(x) \wedge Frau(x)))$$
\end{itemize}

\subsection{Rechenregeln}
\cparagraph{Definition}
\begin{itemize}
\item $\neg \forall x \, F \equiv \exists x\, \neg F$
\item $\neg \exists x \, F \equiv \forall x\, \neg F$
\item Für $Q \in \{\forall, \exists\}$ und $\circ \in \{ \wedge, \vee\}$ gilt:\\
$(Q x F ) \circ G \equiv Qx\,(F\circ G)$, falls $x$ in $G$ nicht frei vorkommt.
\item $\forall x \, F \wedge \forall x \, G \equiv \forall x (F \wedge G)$
\item $\exists x \, F \vee \exists x \, G \equiv \exists x \, (F \vee G)$
\item $\forall x \, \forall y \, F \equiv \forall y \, \forall x \, F$
\item $\exists x \, \exists y \, F \equiv \exists y \, \exists x \, F$
\end{itemize}

\cparagraph{Beispiel}
$(\exists x\,Katze(x)) \vee Vogel(y)\equiv \exists x \, (Katze (x) \vee Vogel(y))$

\subsection{Pränexform}
\cparagraph{Definition} Eine Aussage $F$ ist in \emph{bereinigter Pränexform}, wenn
$$F= Q_1 y_1 \,\dots\, Q_ny_nG$$
wobei $Q_i\in \{\forall, \exists\}$, $G$ keine Quantoren enthält und $y_1, \,\dots\,,\, y_n$ paarweise verschieden sind.\bigskip\\
Für jede Aussage $F$ gibt es eine äquivalente Formel im bereinigter Präexform.
\cparagraph{Beispiel}
\begin{align*}
&\forall x \, (land(x) \to \exists y \, hauptstadt(y,x))\\
\equiv \;&\forall x\, (\neg land(x) \vee \exists y\, hauptstadt(y,x))\\
\equiv \;&\forall x\, (\exists y \, (\neg land(x) \vee hauptstadt(y,x)))\\
\equiv \;&\forall x\, \exists y \, (land (x) \to hauptstadt(y,x))
\end{align*}
\cparagraph{Beispiel}
Nicht jeder Vogel fliegt:
\begin{align*}
&\neg \forall x \, (Vogel (x) \to fliegt(x))\\
\equiv \, & \forall x \, (\neg Vogel(x) \vee fliegt(x))\\
\equiv \, & \exists x \, \neg (\neg Vogel(x) \vee fliegt(x))\\
\equiv \, & \exists x \, (Vogel(x) \wedge \neg fliegt(x))
\end{align*}
(Nicht jeder Vogel fliegt. Durch Umformung erfährt man: Es muss einen Vogel geben, der nicht fliegt)\\
Achtung: Diese Formel sollte nicht wieder in die Implikation umgeformt werden, da sie sonst auch für bspw. einen Stuhl gültig wäre, was nicht gewollt ist: $\exists x (\neg Vogel(x) \to fliegt(x))$ (Stuhl ist kein Vogel, also fliegt er).
\subsection{Hornformeln}
\cparagraph{Definition}
\begin{itemize}
\item Ein \emph{Literal} ist eine Atomformel oder eine negierte Atomformel. Eine nicht negierte Atomformel nennt man \emph{positiv} (auch: Fakt), eine negierte \emph{negativ}. 
\item Eine Formel heißt \emph{Hornklausel}, wenn sie eine $\vee$-Verknüpfung von Literalen ist, von denen höchstens eins positiv ist.
\item Eine Hornformel ist eine $\wedge$-Verknüpfung von Hornklauseln.
\end{itemize}

\cparagraph{Beispiel}
\begin{itemize}
\item Literale: $P(x,y),\; \neg P(x,y), \; \neg Q(z *x+y)$
\item Hornklauseln: $\neg P(x,y) \vee Q(x),\; \neg P(x,y) \vee \neg Q (x), \; Q(x)$
\end{itemize}

\cparagraph{Wichtiger Spezialfall} Hornklauseln mit mindestens zwei Literalen und genau einem positiven Literal. Diese lassen sich als Folgerung darstellen, da
\begin{align*}
&\neg A_1 \vee \dots \vee A_{n-1}\vee A_n\\
\equiv\;& \neg (A_1 \wedge \dots \wedge A_{n-1})\vee A_n\\
\equiv\; & A_1 \wedge \dots \wedge A_{n-1} \to A_n
\end{align*}
Mit Hilfe von Hornklauseln lassen sich Regeln formulieren, z.B. $\forall x \, (Katz(x) \to Haustier(x))$ (die Klammer lässt sich, wie im Spezialfall gesehen, umformen [dieses Mal von der Folgerung zurück), ist sie eine Hornformel]).\\
Wenn diese Regeln mit $\wedge$ verknüpft werden, können diese zusammen mit Fakten (einzelnes positives Literal) als Wissensbasis betrachtet werden.\\
Die Wissenbasis ist dann eine Hornformel.

\cparagraph{Beispiel}\parskp
$katze(feli) \wedge \forall x \, (katze(x) \to haustier(x))$ ist eine Hornformel, die als einfache Wissensbasis betrachtet werden kann. Daraus lässt sich die Aussage dass Feli ein Haustier ist schließen, obwohl es nicht explizit dort steht. Das ist die Aufgabe eines Expertensystems.\\
Grundlegender Aufbau eines Expertensystems:
\begin{center}
\begin{tikzpicture}[scale=.7]
\draw  (-3,3.5) rectangle (-0.5,1.5) node[pos=.5, align=center]{Wissens-\\basis};
\draw  (1,3.5) rectangle (3.5,1.5) node[pos=.5, align=center]{Inferenz-\\system};
\draw [latex-latex] (1,2.5) -- (-0.5,2.5);
\draw [-latex] (5,3) node[right]{Anfrage} -- (3.5,3);
\draw [-latex] (3.5,2) -- (5,2) node[right]{Antwort};
\end{tikzpicture}
\end{center}

\chapter{Prolog-Programmierung}
Prolog: Programming in Logic\\
Prolog-Programme sind im Wesentlichen Hornformeln, bei denen alle Variablen allquantisiert sind. Jede Klausel ist dabei eine Zeile eines Prolog-Programmes.

\cparagraph{Beispiel}$ $
\begin{lstlisting}[language=Prolog]
katze(feli).
haustier(X):-katze(X).
\end{lstlisting}
Dieses Prolog-Programm stellt die Hornformel $katze(feli) \wedge \forall x \, (katze(x) \to haustier(x))$ dar.\\
Anfrage an den Prolog-Interpreter:
\begin{lstlisting}[language=Prolog]
?-haustier(reni).
true.
\end{lstlisting}

\section{Syntax}
Prädikat: Wort in Kleinbuchstaben.\\
Konstante: Wort in Kleinbuchstaben.\\
(Bestimmung ob Prädikat oder Konstante erfolgt über die Position des Wortes: Konstanten treten immer in der Klammer von Prädikaten auf)\\
Variable: Wort, das mit einem Großbuchstaben beginnt.\\
"$\leftarrow$": ":-" (wird impliziert von).\\
"$\wedge$": ",".\\
Jede Klausel wird mit "." abgeschlossen. Alle Klauseln sind implizit mit $\wedge$ verknüpft.













\chapter{Prolog-Programmierung}
Prolog: Programming in Logic\\
Prolog-Programme sind im Wesentlichen Hornformeln, bei denen alle Variablen allquantisiert sind. Jede Klausel ist dabei eine Zeile eines Prolog-Programmes.

\cparagraph{Beispiel}$ $
\begin{lstlisting}[language=Prolog]
katze(reni).
haustier(X):-katze(X).
\end{lstlisting}
Dieses Prolog-Programm stellt die Hornformel $katze(reni) \wedge \forall x \, (katze(x) \to haustier(x))$ dar.\\
Anfrage an den Prolog-Interpreter:
\begin{lstlisting}[language=Prolog]
?-haustier(reni).
true.
\end{lstlisting}

\section{Syntax}
Prädikat: Wort in Kleinbuchstaben.\\
Konstante: Wort in Kleinbuchstaben.\\
(Bestimmung ob Prädikat oder Konstante erfolgt über die Position des Wortes: Konstanten treten immer in der Klammer von Prädikaten auf)\\
Variable: Wort, das mit einem Großbuchstaben beginnt.\\
„ $\leftarrow$ “: „ \lstinline$:-$ “ (wird impliziert von).\\
„ $\wedge$ “: „ \lstinline$,$ “.\\
Jede Klausel wird mit „ \lstinline$.$ “ abgeschlossen. Alle Klauseln sind implizit mit $\wedge$ verknüpft.

\subsection*{$\vee$-Verknüpfung}
\lecdate{03.04.2017}
Die Formel $A\vee B \to C$ ist wegen $A\vee B \to C \equiv (\neg A \wedge \neg B ) \vee C$ keine Hornklausel. Da jedoch $(\neg A \wedge \neg B ) \vee C \equiv (\neg A \vee C) \wedge (\neg B \vee C) \equiv (A\to C) \wedge (B \to C)$, lässt sich $A \vee B \to C$ als Hornformel darstellen.\\
Daher gibt es in Prolog den $\vee$-Operator „ \lstinline$;$ “.
\cparagraph{Beispiel} Die Formel $\forall x \; (katze(x) \vee kater(x) \to tier(x) )$ kann in Prolog als Hornformel dargestellt werden durch:
\begin{lstlisting}[language=Prolog]
tier(X) :- katze(X).
tier(X) :- kater(X).
\end{lstlisting}
Die skann mit dem $\vee$-Operator abgekürzt werden durch \lstinline$tier(X) :- katze(X) ; kater(X).$\\
Die Variable \lstinline$X$ ist implizit allquantisiert (alle Variablen sind in Prolog allquantifiziert).

\subsection{Unterschied zwischen Prolog und imperativen Programmiersprachen}
Eine Variable kann sowohl Eingabe als auch Ausgabe sein. In einer imperativen Programmiersprache lässt sich nachbilden, dass eine Variable entweder Eingabe oder Ausgabe ist.\\
Prolog:
\begin{lstlisting}[language=Prolog]
katze(reni).
katze(mimi).
kater(momo).
tier(X) :- katze(X) ; kater(X).
\end{lstlisting}
Nachbildung in Java:
\begin{lstlisting}[language=Java]
Boolean katze(String x) {
	return x.equals("reni") || x.equals("mimi");
}

Boolean katze(String x) {
	return x.equals("momo");
}

Boolean tier(String x) {
	return katze(x) || kater(x);
}
\end{lstlisting}
Damit können wir die Anfrage für \lstinline$tier("momo")$ stellen, aber keine Anfrage \lstinline$tier(X)$ wie in Prolog, die alle \lstinline$X$ mit \lstinline$tier(X)$ wahr liefert.

\subsection{Existenzquantoren in Prolog durch Skolemisierung}
In Prolog sind alle vorkommenden Variablen allquantisiert. Daher können existenzquantisierte Variablen nicht unmittelbar dargestellt werden. Existenzquantoren können jedoch durch Skolemisierung eliminiert werden. Einfacher Spezialfall: Einführung einer Skolemkonstante.\\
Eine Formel der Form
$$\exists x\; P(X)$$
kann erfüllbarkeitsäquivalent umgeformt werden zu
$$P(a)$$
wobei $a$ eine noch nicht verwendete Konstante ist (Skolemkonstante).\\
So wird die Abfrage $\exists x\; katze(X)$ beispielsweise durch \lstinline$katze(reni)$ umgesetzt.

\section{Auswertstrategie in Prolog}
Die Struktur  eines Prolog-Programmes lässt sich durch einen Und-Oder-Baum darstellen:\bigskip\\
Und-Verknüpfung: \lstinline$a(X) :- b(X) , c(X) , d(X).$\\
\begin{tabular}{L{.3} L{.7}}
\begin{tikzpicture}
\node (v2) at (-0.5,1) {$a$};
\node (v1) at (-1.5,-0.5) {$b$};
\node (v4) at (-0.5,-0.5) {$c$};
\node (v3) at (0.5,-0.5) {$d$};
\draw (v1) -- (v2) -- (v3);
\draw (v4) -- (v2);
\def\wa{-57};\def\wb{-124}\def\rd{.7cm};
\draw[shift=(\wa : \rd)] (-.5,1) arc (\wa:\wb:\rd);
\end{tikzpicture} & \mpb[.6]
Für die Anfrage \lstinline$a(X)$ werden die Zeilziele \lstinline$b(X)$, \lstinline$c(X)$, \lstinline$d(X)$ erzeugt, die alle wahr sein müssen.
\mpe
\end{tabular}\\
Oder-Verknüfung: \lstinline$a(X) :- b(X) ; c(X) ; d(X).$\\
\begin{tabular}{L{.3} L{.7}}
\begin{tikzpicture}
\node (v2) at (-0.5,1) {$a$};
\node (v1) at (-1.5,-0.5) {$b$};
\node (v4) at (-0.5,-0.5) {$c$};
\node (v3) at (0.5,-0.5) {$d$};
\draw (v1) -- (v2) -- (v3);
\draw (v4) -- (v2);
\end{tikzpicture} & \mpb[.6]
Hier muss eines der Teilziele wahr sein.
\mpe
\end{tabular}

\cparagraph{Beispiel} Programm von oben:
\begin{center}
\begin{tikzpicture}
\node (v2) at (-0.5,1) {\lstinline!tier(X)!};
\node (v1) at (-1.5,-0.5) {\lstinline!kater(X)!};
\node (v3) at (0.5,-0.5) {\lstinline!katze(X)!};
\draw (v1) -- (v2) -- (v3);
\end{tikzpicture}
\end{center}
Bei der Anfrage \lstinline$tier(X)$ wird die Variable \lstinline$X$ ersetzt durch die Konstante \lstinline$reni$ und die Variablen auf tieferen Ebenen werden entsprechend ersetzt.\\
Schreibweise: \lstinline$X/reni$ (\lstinline$X$ wird unifiziert mit \lstinline$reni$)
\begin{center}
\begin{tikzpicture}[scale=1.7]
\node (v2) at (-0.5,1) {\lstinline!tier(reni)! ~~ wahr};
\node[align=center] (v1) at (-1.5,-0.5) {\lstinline!kater(reni)!\\ falsch};
\node[align=center] (v3) at (0.5,-0.5) {\lstinline!katze(reni)!\\ wahr};
\draw (v1) -- node[left]{\lstinline!X/reni!} (v2) -- node[right]{\lstinline!X/reni!} (v3);
\end{tikzpicture}
\end{center}
Allgemein kann ein Und-Oder-Baum aus beliebig vielen Und- oder Oder-Verknüpfungen bestehen.\\
Dieser Baum muss systematisch durchsucht werden, um einen Wahrheitswert für die Wurzel zu bestimmen.\\
Mögliche Strategien: Tiefensuche, Breitensuche.
\subsection{Tiefensuche}
Problem bei der Tiefensuche:
\cparagraph{Beispiel} $ $
\begin{lstlisting}[language=Prolog]
zug(X,Y) :- zug (Y,X).
zug(1,2).
\end{lstlisting}
Anfrage \lstinline$zug(2,1)$:
\begin{center}
\begin{tikzpicture}
\node (v1) at (0,3.5) {\lstinline!zug(2,1)!};
\node (v2) at (0,2) {\lstinline!zug(1,2)!};
\node (v3) at (0,0.5) {\lstinline!zug(2,1)!};
\draw (v1) -- (v2) -- (v3);
\end{tikzpicture}
\end{center}
Der Und-Oder-Baum ist unendlich und eine Tiefensuche, die den ersten Zweig (links) verfolgt, findet keine Lösung.\\
Das gleiche Problem tritt bei einer entsprechenden Implementierung in Java auf:
\begin{lstlisting}[language=Java]
Boolean zug(int x, int y) {
	return zug(y,x) || (x==1 && y==2);
}
\end{lstlisting}
\subsection{Breitensuche}
Die Breitensuche würde hier eine Lösung finden:
\begin{center}
\begin{tikzpicture}
\node (v0) at (0,1.5) {\lstinline!zug(2,1)!};
\node (v3) at (-1.5,0) {\lstinline!zug(1,2)!};
\node (v2) at (-3,-1.5) {\lstinline!zug(2,1)!};
\node [align=center] (v4) at (0,-1.5) {\lstinline!zug(1,2)!\\ wahr};
\node (v1) at (-3,-3) {…};
\draw (v1) -- (v2) -- (v3) -- (v0);
\draw (v3) -- (v4);
\end{tikzpicture}
\end{center}
Das heißt, die Breitensuche liefert eine Lösung, die Tiefensuche unter Umständen nicht.\bigskip\\
Lösung: Die Abbruchbedingung der Rekursion muss die erste Klausel sein damit die Rekursion nicht endlos läuft bzw. die Tiefensuche die Abzweigung nimmt, die zu einer Lösung führt:
\begin{lstlisting}[language=Prolog]
zug(1,2).
zug(X,Y) :- zug(Y,X).
\end{lstlisting}


\subsection{Unifizierbarkeit}
\cparagraph{Definition} Zwei Atome $P,Q$ heißen \emph{unifizierbar}, wenn es eine Ersetzung der in $P,Q$ vorkommenden Variablen gibt, so dass damit $P\equiv Q$.

\cparagraph{Beispiel} Mit $a,\; b,\; c$ … Konstanten, $x,\; y ,\; z$ … Variablen.
\begin{itemize}
\item $P(x,a),\; P(a,a)$ sind unifizierbar durch $x/a$.
\item $Q(a,x,y),\; Q(z,b,c)$ sind unifizierbar durch $z/a,\; x/b,\; y/c$.
\item $P(x,x)\; P(a,b)$ sind nicht unifizierbar.
\end{itemize}

$\to$ Auswertstrategie von Prolog:\\
Für ein Ziel $P$ wird das Prolog-Programm von oben nach unten durchsucht, bis eine linke Seite einer Klausel (das, was hinter der Implikation steht bzw. in Prolog vor \lstinline$:-$) mit $P$ unifiziert. Indem die rechte Seite der Klausel ebenso ersetzt wird, werden neue Teilziele erzeugt.

\cparagraph{Beispiel} $ $
\begin{lstlisting}[language=Prolog]
zug(1,2).							% (1)
zug(X,Y):- zug(Y,X).	% (2)
\end{lstlisting}
Für die Anfrage \lstinline$zug(2,1)$ wird folgender Suchbaum erzeugt:
\begin{center}
\begin{tikzpicture}
\node (v1) at (0,3.5) {\lstinline!zug(2,1)!};
\node (v2) at (0,2) {\lstinline!zug(1,2)!};
\node[below] at (v2) {wahr (wegen (1))};
\draw (v1) -- node[right, pos=.5]{\lstinline|X/2, Y/1| (in (2))} (v2);
\end{tikzpicture}
\end{center}
Wenn beispielsweise 
\begin{lstlisting}[language=Prolog]
q(X,Y):- (a,b).
r(Y):- (b).
p(X):- q(X,Y) , r(Y).
\end{lstlisting}
dann:
\begin{center}
\begin{tikzpicture}
\node (v1) at (-1,2.5) {\lstinline|p(a)|};
\node (v2) at (-1,1) {\lstinline|q(a,Y)|};
\node (v3) at (-2,-0.5) {\lstinline|q(a,b)|};
\node (v4) at (0,-0.5) {\lstinline|r(b)|};
\node[below] at (v3) {wahr};
\node[below] at (v4) {wahr};
\draw (v1) -- node[right]{\lstinline|X/a|} (v2) -- node[left]{\lstinline|Y/b|} (v3);
\draw (v1) -- node[right]{\lstinline|Y/b|} (v4);
\def\wa{-56};\def\wb{-124}\def\rd{.7cm};
\draw[shift=(\wa : \rd)] (-1,2.5) arc (\wa:\wb:\rd); %%%TODO
\end{tikzpicture}
\end{center}

\subsection{Backtracking}
Wenn das neue Teilziel \lstinline|false| liefert, wird ein Backtracking ausgeführt und eine andere Alternative gesucht.

\cparagraph{Beispiel}$ $
\begin{lstlisting}[language=Prolog]
q(X):- p(X).	% (1)
q(X):- r(X).	% (2)
p(a).
r(b).
\end{lstlisting}
Anfrage: \lstinline$q(b)$\\
Suchbaum:
\begin{center}
\begin{tikzpicture}
\node [align=center] (v1) at (-1,2.5) {\lstinline|q(b)|\\wahr};
\node [align=center] (v2) at (-2.5,0) {\lstinline|p(b)|\\falsch};
\node [align=center] (v3) at (.5,0) {\lstinline|r(b)|\\wahr};
\draw (v1) -- node[left,pos=.5]{\lstinline|X/b| in (1)} (v2);
\draw (v1) -- node[right, pos=.5]{\lstinline|X/b| in (2)} (v3);
\end{tikzpicture}
\end{center}
Bei einer Und-Verknüpfung im Suchbaum wird zuerst der linke Teilbaum durchsucht, dann die rechten Teilbäume, wobei jeweils die gleichen Ersetzungen von Variablen vorgenommen werden.
\cparagraph{Beispiel} $ $
\begin{lstlisting}[language=Prolog]
p(X):- q(X), r(X).
\end{lstlisting}
Anfrage: \lstinline$p(X)$
Suchbaum:
\begin{center}
\begin{tikzpicture}
\node [align=center] (v1) at (-1,2.5) {\lstinline|p(a)|\\wahr};
\node [align=center] (v2) at (-2.5,0) {\lstinline|q(a)|\\falsch};
\node [align=center] (v3) at (.5,0) {\lstinline|r(a)|\\wahr};
\draw (v1) -- node[left,pos=.5]{\lstinline|X/a|} (v2);
\draw (v1) -- node[right, pos=.5]{\lstinline|X/a|} (v3);
\def\wa{-59};\def\wb{-121}\def\rd{1cm};
\draw[shift=(\wa : \rd)] (-1,2.5) arc (\wa:\wb:\rd);
\end{tikzpicture}
\end{center}
Die Auswertstrategie von oben nach unten lässt sich nutzen, um ein \lstinline$if-else$ zu implementieren.
\cparagraph{Beispiel}$ $
\begin{lstlisting}[language=Prolog]
strassenzustand(regen,nass).
strassenzustand(X,trocken).
\end{lstlisting}
Anfrage 1: 
\begin{center}
\begin{tikzpicture}
\node [align=center] (v1) at (-1,2.5) {\lstinline|strassenszustand(regen,Y)|};
\node [align=center] (v2) at (-1,0) {\lstinline|strassenzustand(regen,nass)|\\ falsch};
\draw (v1) -- node[right,pos=.5]{\lstinline|Y/nass|} (v2);
\end{tikzpicture}
\end{center}
Anfrage 2:
\begin{center}
\begin{tikzpicture}
\node [align=center] (v1) at (-1,2.5) {\lstinline|strassenszustand(sonne,Y)|};
\node [align=center] (v2) at (-1,0) {\lstinline|strassenzustand(sonne,tocken)|\\ wahr};
\draw (v1) -- node[right,pos=.5]{\lstinline|X/sonne, Y/trocken|} (v2);
\end{tikzpicture}
\end{center}
Um zu verhindern, dass auch \lstinline$strassenzustand(regen, trocken)$ zu wahr auswertet, kann die zweite Zeile zu \lstinline$strassenzustand(X,trocken):- X \== regen$ geändert werden.

\section{Darstellung von Relationen in Prolog}
Eine endliche Relation kann durch Aufzählung dargestellt werden. Reflexive, transitive und symmetrische Relationen werden durch entsprechende Regeln dargestellt.
\begin{itemize}
\item Wenn \lstinline$p(X,Y)$ eine Relation ist, kann deren symmetrische Hülle dargestellt werden durch \lstinline$q(X,Y):- p(X,Y); p(Y,X).$
\item Die reflexive Hülle von \lstinline$p$ wird erzeugt, indem die Regel \lstinline$p(X,X)$ hinzugefügt wird.
\item Naiver Versuch Transitivität:
\begin{lstlisting}[language=Prolog]
p(X,Y).
q(X,Z):- p(X,Z); p(X,Y), p(Y,Z).\\
% Achtung: Funktioniert nur für Wege der Länge 2!
\end{lstlisting}
$\to$ Berechnung der transitiven Hülle einer Relation.
\end{itemize}

\subsection{Transitive Hülle}
Sei \lstinline$weg/2$ eine binäre Relation, von der wir die transitive Hülle  berechnen wollen. Es gilt:\\
Es gibt einen Weg von $X$ nach $Y$ genau dann wenn:
\begin{itemize}
\item $X=Y$ oder
\item es gibt einen Knoten $z$ und einen Weg von $X$ nach $Z$ und von $Z$ nach $Y$.
\end{itemize}
Wir müssen daher die reflexive und transitive Hülle von \lstinline$weg$ berechnen.\\
Als Formel:
\begin{align*}
&\forall X\; weg(X,X) \wedge\\
&\forall X \;\forall Y \; ((\exists Z\; weg(X,Z) \wedge weg(Z,Y)) \to weg(X,Y))
\end{align*}
Um den Existenzquantor zu beseitigen formen wir die zweite Klausel um:
\begin{align*}
&\forall X\; \forall Y\; (\forall Z\; \neg weg(X,Z) \vee \neg weg(Z,Y) \vee weg(X,Y))\\
\equiv & \forall X\; \forall Y\; \forall Z\;(\neg weg(X,Z) \vee \neg weg(Z,Y) \vee weg(X,Y))\\
\equiv& \forall X\; \forall Y\; \forall Z\;( weg(X,Z) \wedge weg(Z,Y) \to weg(X,Y))
\end{align*}
In Prolog bspw.:
\begin{lstlisting}[language=Prolog]
weg(a,b).
weg(b,c).
weg(c,d).
weg(X,X).
weg(X,Y):- weg(X,Z), weg(Z,Y).
\end{lstlisting}
Anfrage 1:
\begin{center}
\begin{tikzpicture}
\node [align=center] (v1) at (-1,2.5) {\lstinline|weg(a,c)|};
\node [align=center] (v2) at (-2.5,0) {\lstinline|weg(a,Z)|};
\node [align=center] (v3) at (-2.5,-2) {\lstinline|weg(a,b)|\\wahr};
\node [align=center] (v4) at (.5,0) {\lstinline|weg(Z,c)|};
\node [align=center] (v5) at (.5,-2) {\lstinline|weg(b,c)|\\wahr};
\draw (v1) -- node[left,pos=.5]{\lstinline|X/a|} (v2);
\draw (v2) -- node[left, pos=.5]{\lstinline|Z/b|} (v3);
\draw (v1) -- node[right,pos=.5]{\lstinline|Y/c|} (v4);
\draw (v4) -- node[right, pos=.5]{\lstinline|Z/b|} (v5);
\def\wa{-59};\def\wb{-121}\def\rd{1cm};
\draw[shift=(\wa : \rd)] (-1,2.5) arc (\wa:\wb:\rd);
\end{tikzpicture}
\end{center}
Anfrage 2:
\begin{center}
\begin{tikzpicture}
\node [align=center] (v1) at (-1,2.5) {\lstinline|weg(a,d)|};
\node [align=center] (v2) at (-2.5,0) {\lstinline|weg(a,Z)|};
\node [align=center] (v3) at (-2.5,-2) {\lstinline|weg(a,b)|\\wahr};
\node [align=center] (v4) at (.5,0) {\lstinline|weg(Z,c)|};
\node [align=center] (v5) at (.5,-2) {\lstinline|weg(b,c)|};
\node [align=center] (v6) at (-1,-4.5) {\lstinline|weg(b,Z)|};
\node [align=center] (v7) at (-1,-6.5) {\lstinline|weg(b,c)|\\wahr};
\node [align=center] (v8) at (2,-4.5) {\lstinline|weg(Z,d)|};
\node [align=center] (v9) at (2,-6.5) {\lstinline|weg(c,d)|\\wahr};
\draw (v1) -- node[left,pos=.5]{\lstinline|X/a|} (v2);
\draw (v2) -- node[left, pos=.5]{\lstinline|Z/b|} (v3);
\draw (v1) -- node[right,pos=.5]{\lstinline|Y/d|} (v4);
\draw (v4) -- node[right, pos=.5]{\lstinline|Z/b|} (v5);
\draw (v5) -- node[left, pos=.5]{\lstinline|X/b|} (v6);
\draw (v6) -- node[left, pos=.5]{\lstinline|Z/c|} (v7);
\draw (v5) -- node[right, pos=.5]{\lstinline|Y/d|} (v8);
\draw (v8) -- node[right, pos=.5]{\lstinline|Z/c|} (v9);
\def\wa{-59};\def\wb{-121}\def\rd{1cm};
\draw[shift=(\wa : \rd)] (-1,2.5) arc (\wa:\wb:\rd);
\draw[shift=(\wa : \rd)] (.5,-2) arc (\wa:\wb:\rd);
\end{tikzpicture}
\end{center}

























%\newpage
%\printbibliography

\end{document}